\documentclass[a4paper,12pt]{article}
\usepackage[T2A]{fontenc}
\usepackage[utf8]{inputenc}
\usepackage[russian]{babel}
\usepackage{amsmath,amssymb,mathtools}
\renewcommand{\familydefault}{\sfdefault}
\usepackage{icomma}
\usepackage[margin=18mm]{geometry}
\usepackage{microtype}
\usepackage{xcolor}
\usepackage{tikz}
\usetikzlibrary{arrows.meta,positioning,shapes.geometric,calc}
\usepackage{tabularx,booktabs,longtable}
\usepackage{enumitem}
\usepackage{parskip}
\usepackage{fancyhdr}
\usepackage{setspace}
\usepackage[hidelinks]{hyperref}

\definecolor{accent}{HTML}{276B6E}
\definecolor{darkaccent}{HTML}{17494B}
\definecolor{textgray}{HTML}{2E3335}
\definecolor{softgray}{HTML}{EEF2F1}

\color{textgray}
\onehalfspacing
\setlength{\parindent}{0pt}
\setlength{\headheight}{24.5pt}
\setlist[itemize]{leftmargin=1.7em,itemsep=0.35em,topsep=0.4em}
\setlist[enumerate]{leftmargin=2.0em,itemsep=0.65em,topsep=0.4em}

\pagestyle{fancy}
\fancyhf{}
\fancyhead[L]{\small\color{darkaccent}Текстовые задачи и метод подстановки}
\fancyhead[R]{\small\color{darkaccent}Кирилл Зиновьев}
\fancyfoot[C]{\small\thepage}
\renewcommand{\headrulewidth}{0.4pt}
\renewcommand{\headrule}{\hbox to\headwidth{\color{accent}\leaders\hrule height \headrulewidth\hfill}}

\newcommand{\checkitem}{\item[$\square$]}
\newcommand{\key}[1]{\textbf{\color{darkaccent}#1}}
\newcommand{\sepLine}{\vspace{0.35em}{\color{accent}\hrule height 0.6pt}\vspace{0.8em}}

\tikzset{
  startstop/.style={ellipse, draw=accent, line width=0.9pt, minimum width=38mm, minimum height=9mm, align=center, inner sep=3pt},
  process/.style={rounded corners=3pt, rectangle, draw=accent, line width=0.9pt, text width=88mm, minimum height=10mm, align=center, inner sep=4pt},
  decision/.style={diamond, aspect=2.35, draw=accent, line width=0.9pt, text width=36mm, align=center, inner sep=2pt},
  arrow/.style={-{Stealth[length=3mm,width=2mm]}, line width=0.9pt, draw=darkaccent}
}

\begin{document}

% ---------- Title page ----------
\thispagestyle{empty}
\begin{center}
\vspace*{25mm}
{\Large\color{accent}\textbf{ЧЕК-ЛИСТ}}\\[8mm]
{\huge\bfseries Текстовые задачи}\\[3mm]
{\LARGE\bfseries Составление уравнений и метод подстановки}\\[13mm]
{\large Алгебра, 7 класс}\\[4mm]
{\large 3 октября 2026 года}\\[20mm]

\begin{minipage}{0.78\textwidth}
\centering
\large
Главная цель: научиться переводить условие задачи на язык переменных и уравнений, затем заменять одну переменную выражением через другую и получать уравнение с одной неизвестной.
\end{minipage}

\vfill
{\large\textbf{Лёвин Артём Александрович}}\\[2mm]
{\normalsize эксклюзивно для Кирилла Зиновьева}

\vspace{12mm}
\begin{tikzpicture}[x=1cm,y=1cm]
\draw[accent, line width=1.2pt]
  (-6,0) .. controls (-3.7,0.85) and (-1.8,-0.65) .. (0,0)
         .. controls (2.1,0.75) and (4.1,-0.55) .. (6,0.15);
\end{tikzpicture}
\end{center}
\clearpage

% ---------- Main checklist ----------
\section*{1. Что нужно уметь}
\sepLine

\begin{itemize}
\checkitem Внимательно прочитать условие и определить, \key{что именно неизвестно}.
\checkitem Ввести переменные и обязательно подписать их смысл и единицы измерения.
\checkitem Перевести каждую важную фразу условия в математическое равенство.
\checkitem Получить два связанных уравнения с двумя неизвестными.
\checkitem Выбрать уравнение, из которого одна переменная уже выражена через другую или легко выражается.
\checkitem Выполнить \key{подстановку}: заменить переменную равным ей выражением.
\checkitem Решить полученное уравнение с одной неизвестной.
\checkitem Вернуться к связывающему уравнению и найти вторую неизвестную.
\checkitem Проверить, что найденные величины удовлетворяют и сумме, и разности из условия.
\checkitem Записать ответ словами и с нужными единицами.
\end{itemize}

\section*{2. Как переводить слова в формулы}
\sepLine

\renewcommand{\arraystretch}{1.32}
\begin{tabularx}{\linewidth}{@{}>{\raggedright\arraybackslash\bfseries\color{darkaccent}}p{0.28\linewidth} X@{}}
\toprule
Фраза в условии & Математическая запись \\
\midrule
«Всего» & Складываем все части: например, $x+y=392$. \\
«На $a$ больше» & Большая величина равна меньшей плюс $a$: $x=y+a$. \\
«На $a$ меньше» & Меньшая величина равна большей минус $a$: $x=y-a$. Часто удобнее записать $y=x+a$. \\
«Периметр» & Складываем длины всех сторон. Для равнобедренного треугольника с равными сторонами $x$ и основанием $y$: $2x+y=P$. \\
Есть фиксированная часть & Она тоже входит в целое: например, $x+y+703=6940$. \\
\bottomrule
\end{tabularx}

\vspace{1em}
\textbf{Контрольный вопрос.} После составления уравнения спросите себя: \emph{«Все ли части, о которых говорится в условии, я учёл?»}

\section*{3. Подстановка: главная идея}
\sepLine

Если известно, что
\[
y=x-36,
\]
то в любом другом уравнении вместо $y$ можно написать $x-36$.

Например,
\[
\begin{cases}
 x+y=392,\\
 y=x-36.
\end{cases}
\qquad\Longrightarrow\qquad
x+(x-36)=392.
\]

Теперь неизвестная только одна:
\[
2x-36=392,\qquad 2x=428,\qquad x=214.
\]
Возвращаемся к связи $y=x-36$:
\[
y=214-36=178.
\]
\textbf{Проверка:} $214+178=392$, а $214-178=36$.

\section*{4. Типичные ошибки}
\sepLine

\begin{itemize}
\checkitem Перепутать направление фразы «на $a$ больше/меньше».
\checkitem Написать переменные, но не определить, что означает каждая из них.
\checkitem Забыть одну из частей целого: например, наземный участок трассы.
\checkitem Подставить случайное число вместо выражения, связывающего $x$ и $y$.
\checkitem После нахождения одной переменной забыть найти вторую.
\checkitem Не проверить полученный ответ по исходному условию.
\end{itemize}

\clearpage

% ---------- Flowchart ----------
\thispagestyle{plain}
\begin{center}
{\LARGE\bfseries\color{darkaccent}Алгоритм: как решать текстовую задачу методом подстановки}
\end{center}
\vspace{3mm}

\begin{center}
\small
\begin{tikzpicture}[node distance=5mm]
\node (start) [startstop] {Прочитайте условие целиком};
\node (vars) [process, below=of start] {Введите $x$ и $y$. Запишите словами, что означает каждая переменная.};
\node (eqs) [process, below=of vars] {Переведите данные задачи в уравнения: сумма, разность, периметр, фиксированные части.};
\node (complete) [decision, below=7mm of eqs] {Все данные условия учтены?};
\node (fix) [process, right=16mm of complete, text width=43mm] {Вернитесь к условию, рисунку или схеме и добавьте пропущенную связь.};
\node (sub) [process, below=7mm of complete] {Выразите одну переменную через другую и подставьте это выражение во второе уравнение.};
\node (solve) [process, below=of sub] {Решите уравнение с одной неизвестной.};
\node (back) [process, below=of solve] {Подставьте найденное значение обратно и найдите вторую неизвестную.};
\node (check) [decision, below=7mm of back] {Сумма, разность и другие условия выполняются?};
\node (answer) [startstop, below=7mm of check] {Запишите ответ с единицами};

\draw[arrow] (start) -- (vars);
\draw[arrow] (vars) -- (eqs);
\draw[arrow] (eqs) -- (complete);
\draw[arrow] (complete) -- node[left, font=\small]{да} (sub);
\draw[arrow] (complete.east) -- node[above, font=\small]{нет} (fix.west);
\draw[arrow] (fix.north) |- ($(eqs.east)+(0,-2mm)$);
\draw[arrow] (sub) -- (solve);
\draw[arrow] (solve) -- (back);
\draw[arrow] (back) -- (check);
\draw[arrow] (check) -- node[left, font=\small]{да} (answer);
\draw[arrow] (check.west) -- ++(-28mm,0) node[above, font=\small, midway]{нет} |- (eqs.west);
\end{tikzpicture}
\end{center}
\clearpage

% ---------- Worked examples ----------
\section*{5. Два опорных примера}
\sepLine

\subsection*{Пример 1. Дома на двух улицах}
На двух улицах восстановили 19 домов. На Парковой восстановили на 3 дома меньше, чем на Молодёжной.

Пусть $x$ --- число домов на Парковой, $y$ --- на Молодёжной. Тогда
\[
\begin{cases}
 x+y=19,\\
 y=x+3.
\end{cases}
\]
Подстановка:
\[
x+(x+3)=19,\qquad 2x=16,\qquad x=8,
\]
\[
y=8+3=11.
\]
\textbf{Ответ:} 8 домов и 11 домов.

\subsection*{Пример 2. Трасса и два тоннеля}
Общая длина трассы --- $6940$ м. Наземная часть --- $703$ м. Один тоннель на $17$ м длиннее другого.

Пусть $x$ --- длина первого тоннеля, $y$ --- второго. Тогда
\[
\begin{cases}
 x+y+703=6940,\\
 x=y+17.
\end{cases}
\]
Подстановка:
\[
(y+17)+y+703=6940,
\]
\[
2y+720=6940,\qquad 2y=6220,\qquad y=3110.
\]
\[
x=3110+17=3127.
\]
\textbf{Проверка:} $3127+3110+703=6940$.

\vspace{1em}
\textbf{Перед тренировкой.} Для каждой новой задачи сначала подпишите смысл $x$ и $y$, затем найдите фразу про целое и фразу, связывающую неизвестные. После вычислений обязательно проверьте исходные условия.

\clearpage

% ---------- Training ----------
\section*{Тренировочный блок --- Путь А}
\sepLine
\textbf{10 задач.} В каждой задаче сначала определите смысл $x$ и $y$, затем составьте систему и решите её методом подстановки. Решения записывайте полностью.

\begin{enumerate}
\item В двух коробках 54 карандаша. В первой коробке на 8 карандашей больше, чем во второй. Сколько карандашей в каждой коробке?

\item На двух полках 71 книга. На первой полке на 13 книг меньше, чем на второй. Сколько книг на каждой полке?

\item В двух кассах продали 426 билетов. В первой кассе продали на 38 билетов больше, чем во второй. Сколько билетов продали в каждой кассе?

\item На двух участках высадили 47 деревьев. На втором участке высадили на 9 деревьев больше, чем на первом. Сколько деревьев высадили на каждом участке?

\item Периметр равнобедренного треугольника равен $28{,}6$ см. Каждая из равных сторон на $2{,}9$ см больше основания. Найдите длины сторон треугольника.

\item Периметр прямоугольника равен 74 см. Длина прямоугольника на 9 см больше ширины. Найдите длину и ширину.

\item Длина трассы равна 5250 м. Она состоит из двух тоннелей и наземного участка длиной 620 м. Первый тоннель на 110 м длиннее второго. Найдите длину каждого тоннеля.

\item Три участка дороги имеют общую длину 184 км. Средний участок имеет длину 36 км, а первый участок на 12 км длиннее третьего. Найдите длины первого и третьего участков.

\item Периметр равнобедренного треугольника равен $42{,}5$ см. Каждая из равных сторон на $4{,}75$ см больше основания. Найдите длины всех сторон.

\item Общая длина маршрута равна 8420 м. Два открытых участка вместе имеют длину 920 м, а оставшаяся часть маршрута проходит по двум тоннелям. Первый тоннель на 140 м длиннее второго. Найдите длину каждого тоннеля.
\end{enumerate}

\clearpage

% ---------- Answers ----------
\section*{Ответы}
\sepLine

\renewcommand{\arraystretch}{1.4}
\begin{tabularx}{\linewidth}{@{}c X@{}}
\toprule
\textbf{№} & \textbf{Ответ} \\
\midrule
1 & 31 карандаш в первой коробке, 23 --- во второй. \\
2 & 29 книг на первой полке, 42 --- на второй. \\
3 & 232 билета в первой кассе, 194 --- во второй. \\
4 & 19 деревьев на первом участке, 28 --- на втором. \\
5 & Равные стороны по $10{,}5$ см, основание $7{,}6$ см. \\
6 & Длина 23 см, ширина 14 см. \\
7 & Первый тоннель 2370 м, второй 2260 м. \\
8 & Первый участок 80 км, третий 68 км. \\
9 & Равные стороны по $15{,}75$ см, основание 11 см. \\
10 & Первый тоннель 3820 м, второй 3680 м. \\
\bottomrule
\end{tabularx}

\vfill
\begin{center}
\small\color{darkaccent}
Перед сдачей работы проверьте каждую задачу двумя способами: по сумме всех частей и по указанной в условии разности.
\end{center}

\end{document}
